% Generated draft: requires mathematical and source review.
\documentclass[11pt]{article}
\usepackage[margin=25mm]{geometry}
\usepackage{fontspec}
\setmainfont{Latin Modern Roman}
\usepackage{amsmath,amssymb,mathtools,unicode-math}
\setmathfont{Latin Modern Math}
\usepackage{enumitem,needspace,xurl,hyperref}
\setlength{\parindent}{0pt}
\setlength{\parskip}{5pt}
\setlength{\emergencystretch}{3em}
\newcommand{\N}{\mathbb{N}}
\newcommand{\Z}{\mathbb{Z}}
\newcommand{\Q}{\mathbb{Q}}
\newcommand{\R}{\mathbb{R}}
\newcommand{\C}{\mathbb{C}}
\newcommand{\F}{\mathbb{F}}
\newcommand{\A}{\mathbb{A}}
\newcommand{\PP}{\mathbb{P}}
\newcommand{\E}{\mathbb{E}}
\newcommand{\eps}{\varepsilon}
\newcommand{\dd}{\,\mathrm{d}}
\begin{document}
\section*{UnsolvedMath ID: 30003601}
\textbf{Source problem code:} catalogue entry
\subsection{Definitions and mathematical statement}
Use the usual convention that a standard graded algebra is an algebra over a field \(k\):
\[
R=\bigoplus_{n\geq 0}R_n,\qquad R_0=k,\qquad R=k[R_1],\qquad \dim_k R_1<\infty.
\]
The supplied source does not specify the field \(k\) or any restriction on its characteristic; the intended range of ground fields remains unresolved.

Assume that \(R\) is an integral domain and is Cohen--Macaulay, meaning that every local ring \(R_{\mathfrak p}\) has depth equal to its Krull dimension. Set \(d=\dim R\). Define the \(h\)-vector \((h_0,\ldots,h_s)\) by the Hilbert series identity
\[
\sum_{n\geq 0}\dim_k(R_n)t^n
=\frac{h_0+h_1t+\cdots+h_st^s}{(1-t)^d},
\qquad h_s\neq 0.
\]
Let \(X=\operatorname{Proj}(R)\), the projective scheme associated with the grading. The hypothesis that \(X\) is normal means that each local ring \(\mathcal O_{X,x}\), for \(x\in X\), is an integrally closed domain.

\textbf{Question.} Must every such \(R\) with normal \(\operatorname{Proj}(R)\) have a log-concave \(h\)-vector, that is,
\[
h_i^2\geq h_{i-1}h_{i+1}
\qquad\text{for every integer }i\text{ with }1\leq i\leq s-1?
\]
For \(s\leq 1\), this condition is vacuous. This formulation uses the usual field-based meaning of standard graded; it does not resolve the omitted ground-field qualifications. The supplied background concerns polytope normality and does not specify an additional hypothesis or an equivalent formulation of this question.
\subsection{Short English statement}
Does every standard graded Cohen--Macaulay domain whose associated projective scheme is normal have a log-concave Hilbert-series numerator? The source leaves the ground field unspecified. Its supplied open-status label and unrelated polytope literature assessment require independent review.
\subsection{Sources}
\begin{itemize}
\item Local registry: \texttt{data/unsolvedmath/problems.json}, supplied record with permanent UnsolvedMath ID \(30003601\), title ``Log-Concavity of Cohen--Macaulay H-Vectors,'' source problem code \texttt{catalogue entry}, in the Oberwolfach Reports Open Problems collection. No individual proposer is supplied; the proposed year is recorded as 2017.
\item Supplied report citation: \emph{Mini-Workshop: Lattice Polytopes: Methods, Advances, Applications} (2018), \emph{Oberwolfach Reports}, pp.\ 2659--2701, DOI \texttt{10.4171/owr/2017/44}. Supplied source URL: \url{https://doi.org/10.4171/owr/2017/44}. The record describes Oberwolfach Reports as documenting mathematical workshops and their open-problem discussions. No problem-specific page or named contributor is supplied.
\item Supplied literature source: ``Hofscheier--Kasprzyk, smooth polytopes and IDP,'' \url{https://arxiv.org/abs/2512.20680}.
\item Supplied literature source: ``Bruns, questions on normal polytopes,'' \url{https://www.home.uni-osnabrueck.de/wbruns/brunsw/pdf-article/2-quest-WB.pdf}.
\item Supplied literature source: ``original OWR report,'' \url{https://doi.org/10.4171/OWR/2017/44}.
\item Literature assessment dated exactly \texttt{2026-08-21}: this is supplied dated curation, not independently verified. The record labels the problem open. Its assessment states that the existence of a very ample nonnormal simple polytope would give a counterexample to Oda's conjecture, that no such example is known, and that the general smooth-polytope normality conjecture remains open. These are attributed statements from the supplied assessment; their accuracy and relevance to the present \(h\)-vector question are unresolved. The background itself calls this a dated literature triage derived from supplied OWR curation and requests independent verification of time-sensitive status and cited claims.
\end{itemize}
\end{document}
